\subsection{Representaciones tridimensionales}

Zhang et al. \cite{GCA} proponen un modelo generativo tridimensional denominado \textit{Generative Cellular Automata} (GCA), que extiende el paradigma de los autómatas celulares hacia la generación progresiva de formas 3D. A diferencia de enfoques que intentan reconstruir directamente la forma completa, GCA formula el proceso como una cadena de Markov, donde cada paso corresponde a una transición local sobre una rejilla tridimensional. Estas transiciones se definen mediante reglas homogéneas aprendidas, aplicadas iterativamente sobre las celdas ocupadas de una configuración inicial arbitraria.

El modelo se entrena para que estas reglas locales, al repetirse, converjan hacia formas plausibles de la distribución de datos. Para ello, se representa la forma como puntos de superficie en una rejilla 3D, y se emplea una red convolucional dispersa (Sparse CNN) que opera únicamente sobre las regiones activas.

Durante el entrenamiento, se utilizan secuencias ligeramente distintas a las cadenas de muestreo reales, pero que convergen hacia las formas completas del conjunto de datos. Esto permite que el modelo aprenda reglas de transición robustas, capaces de generar formas diversas desde configuraciones parciales. A través de una sola inferencia, el sistema puede muestrear instancias variadas de la distribución aprendida, logrando resultados competitivos en tareas de completado y generación de formas 3D.

Este enfoque representa una contribución relevante al campo de la morfogénesis computacional, al demostrar que las reglas locales auto-organizadas pueden escalar hacia dominios tridimensionales complejos mediante arquitecturas neuronales adaptadas.

Por otro lado, Shyam Sudhakaran et. al. \cite{NCA3D} expanden el trabajo de \cite{mordvintsevNCA} a tres dimensiones utilizando como espacio celular el juego de Minecraft\footnote{Permite crear estructuras estáticas (como arboles) y estructuras dinámicas (como maquinas voladoras)}. El NCA entrenado puede construir estructuras de hasta 3584 bloques con 50 bloques únicos conservando propiedades regenerativas de su contraparte en 2D.

Una entidad en Minecraft se representa como una red tridimensional donde cada elemento es un vector que almacena (1) el tipo de bloque (incrustado en un vector \textit{one-hot}), (2) su estatus (replicando el comportamiento del canal $\alpha$ de \cite{mordvintsevNCA}) y (3) estados ocultos.

Se utilizan convoluciones 3D para permitir la generación y manipulación de estructuras tridimensionales. Además, se sustituye el uso de filtros de Sobel por una red de percepción aprendible, implementada como una capa de convolución con kernel de tamaño 3, \textit{stride} de 1 y canales de salida equivalentes a tres veces el número de canales del estado celular. Esta red permite que el sistema aprenda activamente a percibir su entorno local, adaptando sus filtros durante el entrenamiento y capturando patrones espaciales relevantes de manera dinámica.

Las características extraídas por esta red perceptual se procesan mediante dos capas lineales configurables, también implementadas como convoluciones 3D con kernel de 1 y \textit{stride} de 1, que generan las actualizaciones del estado celular. Estas actualizaciones se aplican a través de una capa de actualización no configurable, que consolida los cambios en el sistema.

Se utiliza una actualización estocástica y un umbral de activación similar a \cite{mordvintsevNCA}. Finalmente, todos los modelos se inicializan con una distribución normal estándar, con media cero y desviación estándar en el rango de $0.0001$ a $0.1$, lo que proporciona un punto de partida estable para el entrenamiento y la evolución del sistema. Este sistema hereda las propiedades regenerativas de su versión 2D.

Horibe et. al. \cite{horibe} implementan el enfoque regenerativo en robots blandos mediante NCA. Cada robot está definido por un arreglo de volúmenes de reposo para cada vóxel, que se activan mediante una señal global sinusoidal, provocando expansión y contracción volumétrica. Esta dinámica permite simular locomoción a través de ciclos de actuación, evaluando el desplazamiento del centro de masa como métrica de fitness. Los robots blandos se componen de los siguientes voxeles:
\begin{itemize}
	\item Rojo: Músculos activos que se contraen y expanden en fase.
	\item Verde: Músculos que se activan en contra-fase respecto a los rojos.
	\item Azul oscuro: Hueso con alta rigidez.
	\item Azul claro: Hueso con baja rigidez.
\end{itemize}

El modelo entrenado demuestra que los robots blandos pueden recuperar parcialmente su morfología original tras sufrir daño, y restaurar su capacidad de movimiento sin intervención externa. Este comportamiento emerge de la dinámica auto-organizada del NCA, que codifica implícitamente la forma objetivo y su capacidad de reparación.

Para evaluar la capacidad de locomoción y regeneración de los robots blandos generados mediante NCA, se emplea el el simulador físico Voxelyze. Los experimentos se realizan en configuraciones 2D y 3D, utilizando rejillas de $7 \times 7$ y $9 \times 9 \times 9$ vóxeles respectivamente, con condiciones de frontera tipo Dirichlet que limitan el crecimiento morfológico a regiones internas ($5 \times 5$ en 2D y $7 \times 7 \times 7$ en 3D).

Este entorno permite validar la capacidad auto-organizada de los NCA para generar morfologías funcionales y resilientes, y proporciona una base cuantitativa para comparar la eficiencia locomotora antes y después de procesos de regeneración. Además, la estructura voxelizada facilita la integración de reglas locales y codificación distribuida, alineándose con marcos bioinspirados de morfogénesis computacional.